\section{Limitaciones del razonamiento}

\subsection{Los LLM se distraen fácilmente con contexto irrelevante}

\begin{itemize}
    \item Al insertar información irrelevante en un problema de matemáticas (como ``Mario's moustache costs \$10''), el LLM se deja engañar.
    \item Agregar la instrucción ``ignore irrelevant context'' mitiga el problema parcialmente, pero su efecto es limitado cuando aumenta la cantidad de contenido irrelevante.
    \item Incluso la acumulación masiva de oraciones irrelevantes sencillas (``the sky is blue'') provoca una caída significativa del desempeño.
\end{itemize}

\begin{warningbox}{Interferencia del contexto}
El LLM, como modelo probabilístico, incorpora todos los tokens de entrada al cálculo de atención. La información irrelevante no solo desperdicia la ventana de contexto, sino que interfiere de manera sustancial con el proceso de razonamiento. Esto exige especial atención en las aplicaciones reales de Agentes: el ruido de los documentos que introduce la generación aumentada por recuperación (RAG) puede, por el contrario, perjudicar la calidad del razonamiento.
\end{warningbox}

\subsection{El LLM no puede autocorregir su razonamiento}

Los experimentos del equipo de Denny Zhou muestran que:

\begin{itemize}
    \item Cuando se le pide al LLM que revise y corrija su propia respuesta, puede corregir una respuesta incorrecta, pero también puede convertir una respuesta correcta en incorrecta.
    \item En bases de referencia como GSM8K y CommonsenseQA, los métodos de self-correction \textbf{no produjeron ninguna mejora neta}; al contrario, empeoraron el desempeño.
    \item Las mejoras de self-correction reportadas en la literatura suelen usar un Oracle (es decir, solo se le pide al modelo que corrija cuando la respuesta ya es incorrecta), pero el modelo mismo no puede determinar si su respuesta es correcta.
\end{itemize}

\begin{importantbox}{La retroalimentación externa es condición para la autocorrección}
La autocorrección del LLM requiere retroalimentación externa de un Oracle (como las pruebas unitarias en tareas de código). El trabajo de Self-Debug ofrece este Oracle de manera natural mediante pruebas unitarias. Depender únicamente del juicio del propio modelo para corregirse no es confiable.
\end{importantbox}

\subsection{El Multi-Agent Debate no supera a Self-Consistency}

\begin{itemize}
    \item Varios Agentes debaten entre sí y llegan a un consenso (multi-agent debate), generando en total $n$ respuestas.
    \item Sin embargo, aplicar simplemente una votación de Self-Consistency sobre $n$ muestras independientes supera de manera consistente al multi-agent debate.
    \item La interacción de información que introduce el proceso de debate no aporta ningún beneficio adicional.
\end{itemize}

\subsection{El orden de las premisas (Premise Order) afecta el razonamiento}

\begin{itemize}
    \item \textbf{Experimento con GSM8K}: al reordenar únicamente las oraciones del problema (sin cambiar el significado), la exactitud cae alrededor de 10 puntos porcentuales.
    \item \textbf{Experimento de razonamiento lógico}: usando razonamiento puramente lógico con reglas de símbolos aleatorios, con solo alterar el orden de las reglas relevantes, todos los modelos de frontera muestran una caída de más de 30 puntos porcentuales en la exactitud.
    \item \textbf{Causa raíz}: el LLM solo puede procesar la entrada de manera secuencial y no puede, como los humanos, saltar libremente entre premisas ni retroceder sobre ellas.
\end{itemize}

\begin{warningbox}{Efecto del orden de las premisas}
El razonamiento del LLM depende en gran medida del orden en que se presentan las premisas. Al diseñar sistemas de Agentes, el orden en que se organiza la información de entrada debe alinearse con el orden que requiere el razonamiento; de lo contrario, se produce una pérdida significativa de desempeño.
\end{warningbox}

\subsection{Resumen de la sección}

El razonamiento de los LLM actuales presenta tres limitaciones principales: se distrae fácilmente con contexto irrelevante, no puede autocorregirse de manera confiable y es sensible al orden de las premisas. Estas limitaciones ofrecen una guía importante para el diseño de sistemas de Agentes: es necesario compensar estas debilidades mediante herramientas externas y flujos de trabajo cuidadosamente diseñados.
